% !PRINTED 423
\weberpage{423}
%<<CONT
it has an odd square prime divisor $p^2$ which does not divide $y^2$, and
we have
\[
(y+4\lambda p)^2 + m x^2 = 4M^2 + 8\lambda y p + 16\lambda^2 p^2,
\]
and this number, if $\lambda$ is not divisible by $p$, is indeed divisible
by $p$, but not by $p^2$, and so cannot be a square. Moreover, $\lambda$
can still be so disposed of that $y+4\lambda p$ has no odd common divisor
with $x$, so that the conditions of § 114 are satisfied. Accordingly, for
every $m$, $j(\omega)$ satisfies an equation $F_n(u,u)$ in which $C=\pm 1$,
and from this it follows, by Bd. II, § 149:

VI. \gesperrt{The class invariants are integral algebraic numbers.} % [sic] The German print has „Klassenvarianten“; the intended term is Klasseninvarianten (class invariants).

\begin{center}
\textbf{\S\,116. \gesperrt{Class Number Relations.}}
\end{center}

If we equate the degree, as it results from what precedes for both parts of
equation (10), § 115, we obtain the formula
\weq{1}{N = hk + h'k' + h''k'' + \cdots,}
\weq{2}{\sum h k = 2\sum^{\partial} \frac{\partial}{e}\varphi(e) + \varphi(\sqrt{n}),}
in which $\varphi(\sqrt{n})=0$ is to be set when $n$ is not a square, and
$N$ is determined by (13), (14) or (15) of the preceding section.
This is the \gesperrt{Kronecker class number relation}, to which we wish to give a
somewhat more convenient form\footnote{The class number relations, of
which the one derived here is only the simplest case, were discovered by
\gesperrt{Kronecker} (Crelles Journal, Bd. 57) and considerably
generalised by \gesperrt{Gierster} (Mathematische Annalen, Bd. 21, 22) and
\gesperrt{Hurwitz} (ibid. Bd. 25).}.

Beside $n$ we consider all the values $n'$ that arise from $n$ by removing
a square factor, form equation (2) for them, and add together all the
results so obtained. Here only the value $n'=1$ is to be excluded, in case
$n$ is a square, because $F_1(u,u)=0$.

Let $\delta^2$, then, be any square factor of $n$ different from $n$
itself, and
\[
n = n'\delta^2.
\]
If we decompose $n'$ into two factors $n' = a'\partial'$ and denote by
$e'$ the greatest common divisor of $\partial'$ and $a'$, then we have
%CONT>>
